\section{From the one-query argument to the full construction}\label{sec:design-walk}

\paragraph{What this section proves and how to read it.}
This section proves nothing and is not a premise of any later
section. Its purpose is to explain, before the formal construction in
Part~V, why each feature of that construction is needed.
\Cref{sec:mechanism} reduced the theorem to: exhibit a signed law
$\nu$ whose nonnegative summands accept six disjoint queries across
$\iota$, dominate it at $\iota$ as in Part~III, and achieve
$\Theta<1$. Each subsection below follows the same pattern:
\emph{Problem} means ``with the naive choice just described,
$\Theta\ge1$ (or no scale-independent bound exists), because [an
explicit number]'', and \emph{Construction} describes the choice
used in Part~V and names the section proving it. The numerical illustrations use the density-linear baseline of
\cref{sec:closing}. The final theorem uses the additional moments,
normalized assignments and joint mean allocation in
\cref{sec:four-routes-closing}; \cref{subsec:walk:glance} connects
these ingredients. Historical comparisons are collected in Part~VI.

Every version of the argument has the same three tasks: verify
finitely many local operator inequalities, sum the resulting prices
over blocks of primes, and check $\Theta<1$. The subsections explain
the local laws that make the first task possible with small
constants, the block summation that makes the second finite, and the
sharper inequalities that make the third true.

The earlier $216$-slot construction used the same general tools.
\Cref{tab:walk:substitutions} records the four structural changes;
\cref{subsec:history:216} gives the historical instances and their
certificates. Those instances are not premises of the present proof.

\begin{table}[htbp]
\centering
\small
\caption{The four structural substitutions.}\label{tab:walk:substitutions}
\begin{tabular}{@{}>{\raggedright\arraybackslash}p{0.375\textwidth}%
>{\raggedright\arraybackslash}p{0.45\textwidth}%
>{\raggedright\arraybackslash}p{0.11\textwidth}@{}}
\toprule
$216$-slot design & final design & explained in\\
\midrule
$V=(\Zsix)^3$, sorted by nine folds; the reserved cut is a chamber
reflection of the third coordinate; exponent $a/(216-\theta)$
& $V=(\Zsix)^2\times\{0,1\}$, sorted by seven folds; the reserved cut is the
layer flip $\iota$; exponent $a/(72-\theta)$ ($a/71$ at $\theta=1$)
& \cref{subsec:walk:queries}\\
\addlinespace
nine coupled families and a singleton recurrence, with the records that do
not accept covered by obstruction cylinders; later one fixed family of
sixteen (at one checkpoint four) amplified edges
& six disjoint one-way queries with a common signed index; records that do
not accept are repaired with their actual marginals, or deleted
& \cref{subsec:walk:queries,subsec:walk:repair}\\
\addlinespace
a comparison law that is reflection positive at the reserved cut
& a comparison law that is a signed combination of nonnegative accepting
laws, not reflection positive, controlled at $\iota$ by a reserved Gram
bracket; symmetric and skew changes priced separately; signed whole-block
fold records, variance and envelope credits, common mean
& \cref{subsec:walk:bracket,subsec:walk:accounting}\\
\addlinespace
the same laws for both classes of primes (a map law, two middle arrays, a
three-depth high range); small primes in the index with generic bounds
& class one: a square-clique prefix; class three: a twelve-label middle law
and a two-layer $72$-slot upper law; small class-three primes in the index
with the exact square-shift contraction
& \cref{subsec:walk:blocks,subsec:walk:bracket,subsec:walk:analytic}\\
\bottomrule
\end{tabular}
\end{table}

\subsection{Put square-difference constraints into the reference law (graph laws)}\label{subsec:walk:graph-laws}

\paragraph{Problem.}
The crude accepting law of \cref{subsec:mech:crude} starts from independent
labels and conditions them on a square difference. At every prime its change
has variation more than one, and the prices of the primes multiply. The
reference knows nothing about squares, so the comparison law has to create
every square relation by itself.

\paragraph{Construction: graph laws.}
The square relations are put into the reference. At a prime $p$ the labels
are drawn from a \emph{graph law} $\lambda_p(G)$ (\cref{def:graphs:labeling}):
the uniform law on the labellings of the vertices of a small oriented graph $G$
by distinct elements of $\Fp$ in which every edge $u\to w$ is a strict-square
constraint, $z_w-z_u\in\Qsq$. The slots are labelled through maps from slots to
vertices; the laws obtained in this way are the \emph{records} of
\cref{subsec:walk:completions}. A record that maps a query to an edge of $G$
already forces that query at the prime and is kept unchanged, so
it needs no replacement and contributes no replacement price. The middle law lives on twelve
vertices, $\Zsix\times\{0,1\}$ (one chamber coordinate and the layer), with all
four layer lifts of the six edges $u\to\rot u$: $24$ directed edges, every vertex of
outdegree two and indegree two (four incident constraints), and the
edge $(0,0)\to(2,1)$ onto which every query projects. The
upper law lives on all $72$ slots, with $720$ directed edges and every vertex of
outdegree ten and indegree ten (twenty incident constraints), and
its root contains every query
(\cref{subsec:push:layer12spectral}, \cref{lem:upper-geometry}).

\paragraph{Why a graph law is almost independent.}
The number of labels available to a vertex whose $t$ neighbours carry distinct
labels is $p2^{-t}$, up to a relative error of order $2^tp^{-1/2}$, by Weil's
bound (\cref{lem:extension}). Iterating this count gives every marginal,
conditional law and density of $\lambda_p(G)$, relative to the independent
law, with explicit relative errors of order $2^tr$, $r=p^{-1/2}$
(\cref{sec:labellings}); the relative form is what makes them usable at
explicit thresholds (\cref{subsec:intro:new}). The law is invariant under the maps
$z\mapsto sz+c$ with $s$ a nonzero square, so its one-slot marginals are
uniform (\cref{lem:marginals}). The errors grow like $2^t$ with the number of
constraints at a vertex, which is one reason why graphs of small degree are
used where the primes are small: the twelve-label law from
$Q_3=19\,951\,586$ on, the $72$-slot law only from $\Pthree=167\cdot2^{40}$ on
(\cref{subsec:walk:blocks}). What remains is that a graph law is not
reflection positive (\cref{subsec:walk:completions}), and that the records
which do not map the queries to edges must be changed
(\cref{subsec:walk:repair}).

\subsection{Graph laws are not reflection positive: adding correction summands}
\label{subsec:walk:completions}

\paragraph{Problem.}
At a chamber cut the two halves of a graph law are coupled through its cross
edges. Relative to the product of its half laws, its normalized cut operator
is a combination of character kernels such as $\chi(y-x)$, each of normalized
size about $p^{-1/2}$, and its spectrum reaches down to about $-gr$, for a
constant $g$ of the graph (\cref{sec:cuts}). So $\lambda_p(G)$ is not reflection
positive, and the signed H\"older inequality (\cref{thm:holder}) does not apply
to it.

\paragraph{Construction: the five-type completion.}
The root law (type L, weight $1$) is completed by records built from the laws
of its half graphs and composed with maps of the chamber coordinate onto
parts of itself: three folds (type D), two matchings (S), six merges (R) and
two constants (C), with the graded weights $1$, $w_{\rm D}r$, $w_{\rm S}r$,
$w_{\rm R}r^2$, $w_{\rm C}r^3$ (\cref{lem:five-type}).
Positivity at a chamber cut is built in steps of one order in $r$ each. The
aligned fold is a copy of the half law placed symmetrically; it is fixed
samplewise by the cut, and it supplies the positive diagonal that cancels the
negative part $-dr$ of the root if $w_{\rm D}\ge d$ (here
$d,f,h$ are the spectral constants of \cref{lem:five-type}; this
$d$ is not the slot count $d=72$). A fold that is aligned at
one cut is not aligned at the others, where it couples two corners of the
half; the matchings pay for this if $w_{\rm S}\ge w_{\rm D}$. The row errors
$fr$, $hr$ of the folds and matchings are paid by the merges if
$w_{\rm R}\ge fw_{\rm D}+hw_{\rm S}$, and the linear terms of the merges by the
constants if $w_{\rm C}w_{\rm S}\ge w_{\rm R}^2$; the constants, of
weight $r^3$, together with the matchings, of weight $r$, dominate
the $r^2$ cross terms of the merges in the Schur complement at the
cut, since the product of the $r^3$ and $r$ weights has the same
$r$-order as the square of the $r^2$ merge weight. The chamber group transports
positivity from one chamber cut to all of them (\cref{lem:slot-array}(i)), and
one-layer copies of the completion records, placed on the reserved diagonal, restore
positivity at $\iota$. (The history of this completion, through
an earlier version on $216$ slots, is in \cref{sec:history}.) The
twelve-label reference of the middle range (thirteen inequalities
per piece, together with the inequality that protects the
$2\times2$ block at $\iota$, \cref{lem:front12:reference}) and
the $25$-state two-layer reference of the upper range ($105$
inequalities per band, \cref{lem:upper-reference-rp}) are its two
instances in the construction.

\paragraph{Reweighting the principal law by a nonnegative function.}
The root itself need not be the graph law. Reweighting it at every prime by a
nonnegative weight built from trace scores---the average
$\frac13\sum_j(1+\ulaw\Hs j)^2$ of the individual squares for the three
opposite base pairs, or $1/100+(1+\ulaw H_{\mathcal D_{39}})^2$ for the
coherent dictionary of $39$ traces---keeps its accepting support, its symmetries and its uniform one-slot
marginals, and lowers the constant $d$ of its chamber cut below $\sqrt{30}$
(\cref{def:positive-root,lem:positive-root-properties,lem:imp:coherent39-law,cert:positive-root-caps}). The
price is a larger constant $k_{\rm L}$ of the layer cut, paid through the
weight $v_{\rm L}$ of the one-layer copy (\cref{prop:root-layer}), and a changed
half marginal, paid in the restorations of the variance credits and of the
mean savings (\cref{lem:record:root-marginal}(iv)). Both cut bounds are proved
in the fixed half charts in which the completion inequalities are stated
(\cref{rem:root-charts}). The gain is only in the negative part
of the chamber-cut operator. Indeed, every law satisfying the
sixteen cross constraints of that cut, with the half marginal and
chamber symmetries of the graph law, has centred norm --- the
operator norm on half-functions of mean zero in the stated half
marginal --- at least about $\sqrt{30}\,r$: with those marginals
no improvement is
possible. A nonnegative reweighting, however, changes the half
marginal, and can thereby raise the least eigenvalue
(\cref{subsec:obs:root}). The construction tilts the root on the
seven pieces from $2^{30}$: averages of individual squares on six
of them, and on the final piece the coherent $39$-trace square
with the constant $1/100$ added. Sharper collision bounds on two
explicit subspaces, an interpolation between the two pair bounds,
and a layer remainder summed without expanding all its monomials
make these parameter values work (\cref{sec:positive-roots}).

\paragraph{Inputs.}
The constants $d$, $f$, $h$ are spectral caps and row bounds of graph laws
(for $d$ on the seven last middle pieces, of the positive roots). The
spectral criterion certifies at one prime that the normalized cut operator of
a graph law has spectrum in $[-gp^{-1/2},1+O(p^{-1})]$, by a character expansion
of its kernel, an extraction of witnesses (\cref{lem:matching:extractionk}) and
a Schur test; monotone majorants carry the bound to all larger primes
(\cref{lem:majorant}). It has two routes: Route~A, through
conditional laws of the opposite half
(\cref{thm:spectral-criterion}), and, at the root of the upper law
below $2^{52}$, Route~B, a direct decomposition into explicit
uniform averages over allowed sets
(\cref{thm:record:direct-star}). At the chamber cut of the uniform middle
root, complete actual stars give the cap directly: the eighteen star features
of the root have an exact signed star matrix of norm $\sqrt{30}$, and a
translation-sector completion, a rectangular factorization of the root and
balanced collisions bring the cap at $Q_3$ below $5.500058$
(\cref{subsec:middle-l20}). The row bounds
come from pinned influence, through Kassabov's and Peyre's inequalities
(\cref{prop:middle-negative}), and from density ratios
(\cref{lem:record:class-three,lem:row-comparison}).

\paragraph{Why second-order character sums are needed.}
On a block of primes (\cref{subsec:walk:blocks}) the negative side of a cut
operator is governed by the worst prime, but the positive side multiplies over
the primes: if the factors have spectra in $[-g_p,1+b_p]$, with
$0\le g_p\le1+b_p$ and $b_p\ge0$, the tensor product has spectrum in
\[
 \Bigl[-\max_pg_p\prod_{q\ne p}(1+b_q),\ \prod_p(1+b_p)\Bigr]
\]
(\cref{lem:spec:tensorcomp}). The density ratios of the
conditional rows --- the conditional distributions of one half
given the labels of the other half, in the row notation of
\cref{sec:spectral} --- multiply over the block in the same way. These quantities must therefore be of
order $1/p$, and a first-order bound, such as Weil's bound for one character,
of order $p^{-1/2}$, is not enough for them. In the $216$-slot design the
first-order bound for a row of its middle completion would have
given a bound of about $35$ for the logarithm of the product of
the row factors over its first middle piece, where the
second-order bound gave $4.3\cdot10^{-3}$ for the same logarithm. In the final design some of the second-order quantities
follow from exact counts at the primes $p\equiv3\pmod4$: the density ratios of
the middle rows lie between $R_2^{-1}$ and
$R_2=(p-3)(p-7)/((p-11)(p-15))=1+O(1/p)$, and between $R_1^{-1}$ and
$R_1=(p-1)/(p-3)$ for one layer (\cref{lem:record:class-three}). The others need
genuinely second-order estimates (\cref{sec:dl}; \cref{subsec:intro:deligne}
says where each enters): the four-cycle kernel, bounded through the character
table of $\mathrm{GL}_2(\Fp)$ and Jacobi sums (\cref{thm:dl:K}); character
sums in two variables and second moments of trace functions with
a moving root $\lambda$ (not the principal law), bounded through
Deligne's theorem and the Katz--Sarnak description of a family
of curves (\cref{thm:dl:H2,lem:moving-root-moments}); and, for the second moment
of the neighbourhood count behind the half variances of the root, a twisted
second moment of Legendre traces, evaluated through a cusp form of weight four
(\cref{thm:ahlgren-ono,lem:record:weighted-moment}).

\subsection{Summing a price at every prime is too large: group primes into blocks}\label{subsec:walk:blocks}

\paragraph{Problem.}
With graph laws and repairs (\cref{subsec:walk:repair}) the comparison law
changes a record of weight of order $r$ by a measure of norm of order $r$, so
the change at one prime costs about $c/p$, for a local constant $c$, and is
priced $(c/p)\,p^{2a}$. The sum $c\sum_pp^{2a-1}$ over all primes is unbounded
(\cref{subsec:mech:crude}), and on a finite range it is large: over the primes
$p\equiv3\pmod4$ of $[Q_3,\Pthree)$ the sum $\sum_p1/p$ alone is about $0.33$,
while the local constants of the middle records are of order ten (twice a
completion weight $w_{\rm D}$ between $5.43$ and $5.51$, times a norm
coefficient of order one).

\paragraph{Construction: blocks.}
Consecutive primes are grouped into blocks $\Blk$ with $\log q_\Blk\le L$,
$q_\Blk=\prod_{p\in\Blk}p$ (\cref{subsec:record:middle_blocks}), and each
block carries one law: the completion,
with $r=b_\Blk^{-1/2}$ for the least prime $b_\Blk$ of $\Blk$, of the tensor
product over its primes of the graph laws (for the root on the seven last
middle pieces, of their positive tilts). A tensor product of operators that
preserve the constants
decomposes into exact-support sectors (\cref{lem:sector}). On a sector with one
active prime, the difference of two tensor products is the difference at that
prime; on a sector with several active primes it is a product of small
factors. So a change made jointly at all primes of a block costs about as much
as at its worst prime, $c/b_\Blk$, and its price is
$c\,b_\Blk^{-1}q_\Blk^{2a}\le c\,e^{2aL}/b_\Blk$: one selection factor $e^{2aL}$
for the whole block. The negative spectrum of a block is likewise that of its
worst prime, times a product $\prod_{p\in\Blk}(1+O(p^{-1}))$ that stays bounded
(\cref{subsec:walk:completions}).

\paragraph{Choosing the block-size bound $L$.}
Near $x$ about $L/\log x$ primes $p\equiv3\pmod4$ fit into one block, so about
$dx/(2L)$ blocks start in $[x,x+dx]$, and heuristically
\[
 \sum_{\Blk}\frac c{b_\Blk}\,q_\Blk^{2a}\approx
 c\,\frac{e^{2aL}}{2L}\,\log\frac{\Pthree}{Q_3}.
\]
The factor $e^{2aL}/L$ is least at $L=1/(2a)$, where the total is
about $c\,a\,e\log(\Pthree/Q_3)$ and grows linearly with $a$. Hence
$a$ is limited by the local constants $c$: every price grows with $a$
through $q_\Blk^{2a}$, and $a$ can be raised only by lowering a local
constant. At the baseline parameters $e^{2aL}\log(\Pthree/Q_3)/(2L)$
is about $0.039$, against about $0.33$ for single primes. The
construction uses $L=570$, close to $1/(2a)\approx562.8$
(\cref{def:record:exact-blocks}); although every block satisfies $\log q_\Blk\le L$,
the local laws are verified at the slightly larger bound
$L_{\rm loc}=585$, hence with the prime-count cap
$\lfloor L_{\rm loc}/\log X\rfloor$, so they hold on every block
(\cref{lem:record:local-domain}). Middle prices computed at the
evaluation exponent $a_0=\PushAnchor$ are converted to $a$ by
$t=e^{2(a-a_0)L}\approx1.000000062$ (\cref{lem:record:repricing}).
Blocks are formed from exact prime counts below an explicit
threshold and from fixed-size prime-counting intervals above it,
controlled by \cref{thm:bmor}.

\paragraph{The ranges.}
A block is priced with the local constants of its least prime, so different
laws are used where different constants are smallest (\cref{tab:ranges}).
Below $Q_3$ the primes $p\equiv3\pmod4$ are not coordinates of the labels but
enter the index of the queries, each at a factor close to one in the conductor
moment $R_I$ (\cref{subsec:walk:analytic}). In $[Q_3,\Pthree)$ the twelve-label
law needs one added edge in each nonaccepting record of degree one, so its
changes are of order $r^2=1/p$ per prime. From $\Pthree=167\cdot2^{40}$ on, the
two-layer law on all $72$ slots contains every query in its root, every one of
its records of degree at most one accepts. Repairs start at degree two
and deletions at degree three, so its operator prices at one prime
are of order $r^3=p^{-3/2}$
(\cref{sec:frontier:upper-oneway}); the Route~B criterion supplies
its root caps from $\Pthree$ to $2^{52}$, with graded root counts and
coefficients interpolated between certified prime thresholds, and
the Route~A criterion those from $2^{52}$ on
(\cref{subsec:upper-bands}). From
$\Pfour=2^{512}$ on, the primes of both classes are single atoms with a fixed
two-layer tail law, at prices $O(p^{-3/2+2a})$ with a convergent sum
(\cref{sec:tail-index}). The primes $p\equiv1\pmod4$ below $2^{512}$ need no
comparison at first order (\cref{subsec:walk:bracket}).

\subsection{One query leaves an error term above $1$: use six disjoint queries}
\label{subsec:walk:queries}

\paragraph{Problem.}
With one query the coupled term $BR_I/2$ of \cref{prop:mech:one-query}
is about $1.45$ at the baseline values, already above $1$ before any
comparison price is added (\cref{subsec:mech:budget}).

\paragraph{Construction: six disjoint queries.}
With $k$ disjoint queries and independent indices the matching term is
$Be^T(R_I/2)^k$ (\cref{subsec:mech:budget}, \cref{thm:comparison}).
But every retained record of the comparison law must accept every query, and
acceptance costs added constraints. The geometry of $V$ decides how many
queries are cheap. The six queries $e_j=((0,j,0),(2,\rot j,1))$, $j\in\Zsix$,
of \cref{def:slot-array} all have the first-coordinate pair $(0,\rot0)=(0,2)$;
they differ only in the second chamber coordinate, on which they act by
$\rot$. The middle laws ignore that coordinate (\cref{lem:slot-array}(v)), so
they see one query, the edge from the label $(0,0)$ to the label $(2,1)$, and
each nonaccepting fold or matching record of the middle law needs one added
edge; the prefix law ignores
the layer as well and sees one coincidence event
(\cref{subsec:front6:construction}). Six is the largest number of disjoint
queries with this projection, since there are only six slots $(0,j,0)$. So six
queries cost the middle range no more than one, while the matching term
carries $(R_I/2)^6<0.0588$; finally $J<0.25177$.

The larger family of twelve queries $((u,j,0),(\rot u,\rot j,1))$,
$u\in\{0,5\}$, has two first-coordinate pairs, $(0,2)$ and $(5,3)$, and each
nonaccepting fold or matching record then needs two added edges instead of
one. The construction was first proved for these twelve queries and then
specialized to the six, and only the six-query form was certified
(row~8 of \cref{tab:history:checkpoints}). The twelve queries survive as the
acceptance set of the middle merge repairs, of the one-way repairs and the
dense EL comparison of the upper range, and of the tail, while the signed fold
records of the middle, the LE mixture and the retained records $G_{01}L$,
$G_{10}L$ of the upper range accept the six queries only
(\cref{subsec:front12:construction,lem:six-query-retention}). Twelve queries
would need more than new edges: joint laws (couplings) of the two
halves forcing both queries of a fold record do exist, but with
the unchanged references the symmetric error of the
fold record is then of order $r$, at least about $r/\sqrt3$, and that of a
matching fork at least about $\sqrt3\,r$ in symmetric plus skew norm
(\cref{subsec:obs:twelve}).

\paragraph{Why $216$ slots became $72$.}
In the $216$-slot design the laws of the lower ranges ignored the third
chamber coordinate: its only role there was to carry the reserved cut. A
two-point layer with one reflection serves as well. The folds
$f_0\circ f_4\circ f_0$ in each of two chamber coordinates, followed by the one
fold of the layer, map $V=(\Zsix)^2\times\{0,1\}$ to a single slot
(\cref{lem:slot-array}(iii)), so the chessboard estimate gives
$\alpha^{72}\le Y$. Since the exponent of the theorem is $a/(|V|-\theta)$
(\cref{thm:comparison}), the same price exponent $a$ gives about three times
the saving: $a/(216-\theta)$ becomes $a/(72-\theta)$, and the density-linear
induction (\cref{subsec:walk:analytic}) then gives $a/71$. The laws of the high
range, which used the third coordinate in the $216$-slot design (a hierarchy of
three depths), became two-layer laws on the $72$ slots. The first $72$-slot
design, which also treated the two classes of primes separately, raised the
exponent from $2.288\cdot10^{-7}$ to $1.992\cdot10^{-6}$.

\subsection{Records that do not force a square difference: replace, delete, or dominate them}
\label{subsec:walk:repair}

\paragraph{Problem.}
The reference law is a positive combination of records. Most of its
mass already forces the queries (in the middle range, the principal
law and the aligned fold). Some records do not: in the middle range
the other two fold records $F_1,F_2$ and the two matching records
$E_0,E_1$ (degree one), and the merges and constants; in the upper
range some records of degree at least two. Bounding a non-forcing
record by its whole mass (a union bound over obstructions) charges
order $r$ for a degree-one record, which is first order and too
large: summed over primes it cannot give $\Theta<1$ (see
\cref{subsec:obs:polynomials} for the abandoned union-bound
polynomials, which play no logical role). Each non-forcing record
must instead be removed by domination, replaced at fixed marginals,
or deleted, according to its type.

\paragraph{Construction: polarization and replacement at fixed marginals.}
Both are proved in \cref{sec:completion-repair}.
\emph{Polarization} (\cref{lem:tr:polarization}): at a chamber cut the two
nonaligned fold records are dominated, in the cut order, by the two matching
records plus a merge correction of one more order in $r$, so they can be
removed without losing positivity at that cut. This is an order of cut
matrices, not of measures; only actual records are subtracted, so the law
stays nonnegative. \emph{Additive repair} (\cref{lem:tr:repair}): a record whose
graph lacks the query edge is replaced by the graph law with the edge added,
reweighted by a factor $1+f(x)+g(y)$ so that its half marginals at the cut are
exactly those of the old record. The reweighting keeps the support, hence
every constraint, and $f,g$ solve two linear equations that are well
conditioned when the conditional operator of the new record is a contraction.
The new and the old record then have the same actual half marginals, so their
difference annihilates the constants on both sides of the cut; its energy is
controlled by its operator norm between the actual half spaces, which is of
order $r$, not by its mass. A record of weight of order $r$ costs
$O(r)\cdot O(r)=O(1/p)$. Records of high degree are simply deleted and charged
by their mass. The copies of records that the completion places on the
reserved diagonal are handled by a \emph{first stage}: polarization at the
chamber cut $\vartheta_0$ removes their nonaccepting fold records, and a
one-stage repair keeps the actual conditional marginals of their matchings
(\cref{lem:record:first-stage}, \cref{sec:middle-frame}); its loss, with the
errors of the two folded matching repairs tensored at their own first cuts
(\cref{cor:two-cut-series}), is $L_0<4.7634\cdot10^{-4}$.

\paragraph{Replacements that keep the old cross edges.}
A one-way repair builds its new record from two half copies and the projected
queries and discards the old cross graph, so that a degree-two record whose old
cross graph is large changes by a kernel of the size of that whole graph. For
such records the upper comparison keeps the old cross graph, adds the missing
required edges, softens the unrequired old edges that meet them, and repairs
the result to the exact old actual marginals. Written in one
common list of actual star feature functions, the old and the new
star matrices then nearly cancel, and the
difference is priced separately in its symmetric and skew parts. For the six
LE records four such softened laws, each repaired to the same marginals, are
mixed; the star matrix of the mixture is the mixture of their matrices, and at
the selected weights its norm is below that of every component
(\cref{def:dense-comparison,prop:dense-comparison,prop:le-mixture}). At $2^{48}$
the symmetric and skew operator norms in the actual half
marginals are about $1.84r$ and $1.83r$ per LE record and $3.20r$
and $2.31r$ per EL record, against about $9.2r$ per LE record and
$9.7r$ per EL record, charged in full to both prices, for one-way repairs of the
same records at $2^{49}$ in the construction of row~11 of
\cref{tab:history:checkpoints}.


\subsection{The forcing law is not reflection positive: domination by $2\times2$ blocks}
\label{subsec:walk:bracket}

\paragraph{Problem.}
In the $216$-slot design the middle laws and all their changes depended only
on the first two chamber coordinates, so every nonnegative law built from them
was fixed samplewise by the reserved cut, and a nonnegative measure fixed
samplewise by a cut is positive there (\cref{lem:commoncut}(i)): the
comparison law stayed reflection positive at the reserved cut. The final
middle law carries the layer: its twelve labels are a
chamber coordinate times the layer, and its query edge runs from layer $0$ to
layer $1$. So $\iota$ acts on its labels, and a change of a record is, at
$\iota$, a signed measure that is not positive. In the first $72$-slot design
the upper laws already carried the layer, and positivity at $\iota$ was kept
by adding to each repair a positive diagonal of the size of its change, the
one-sided bracket of \cref{subsec:history:layers}; this diagonal was paid at
first order.

\paragraph{Construction: domination at the reserved cut, with symmetric and skew parts bounded separately.}
The comparison law $\nu$ need not be positive at any cut, and on the middle
blocks it is not even nonnegative: there it is a finite signed combination of
nonnegative laws that accept the queries (\cref{def:cp4}). It is compared with
the first stage $\lambda$ at $\iota$ alone, through a nonnegative measure
$\Nbr$ with
$\bigl(\begin{smallmatrix}\Nbr&\lambda\\ \lambda&\Nbr\end{smallmatrix}\bigr)
\succeq_\iota0$, that is $-\Nbr\preceq_\iota\lambda\preceq_\iota\Nbr$
(\cref{lem:record:gram}); its energy factor is $B=e^{U_0}+R_{\rm err}<2.33263$
(\cref{lem:bracket-energy}). Every change $\nu_\Blk-\lambda_\Blk$ is split into
its symmetric and skew parts at $\iota$, and each part is priced with its own
positive envelope (\cref{def:cp4}). A nonzero skew part is never positive
semidefinite, since a cut matrix that is $\succeq0$ is symmetric; keeping $\nu$
positive at $\iota$ means adding positive mass of the size of the change,
which is what the first $72$-slot design paid. But a skew part costs nothing at
first order. In the expansion of $Y_\nu-Y_\lambda$ over the blocks, every term
with an odd number of skew factors vanishes, because both halves carry the
same half-function and $\lambda$ is symmetric at $\iota$
(\cref{lem:record:record-expansion}). The reserved cost is therefore
$B(e^S\cosh A-1)$: the symmetric price $S$ enters linearly, the skew price $A$
only through $\cosh A-1\approx A^2/2$. A comparison law that is positive at
$\iota$ is its own bracket (\cref{rem:mech:positive}); the positive
correction of the first $72$-slot design is that case.

\paragraph{Signed combinations over a whole block of primes.}
So the first-order change of an accepting record should be skew. Write $x_i$
and $y_i$ for the labels of the object $i$ of a two-layer fold record on the
layers $0$ and $1$. At $p\equiv3\pmod4$ the number $-1$ is a nonsquare, so the
reserved reflection, which exchanges $x$ and $y$, maps the channel
$\mathsf C=\chi(y_0-x_1)$ of the added accepting edge $x_1\to y_0$ to
$-\mathsf D$, where $\mathsf D=\chi(y_1-x_0)$ is the channel of its mirror
$x_0\to y_1$. Writing $K_{\rm old}$ for the cut kernel of the reference fold
record at $\iota$, relative to the product of its two actual half
laws, adding both edges multiplies this kernel by
$(1+\mathsf C)(1+\mathsf D)$, so the change of kernel is
\[
 K_{\rm old}(1+\mathsf C)(1+\mathsf D)-K_{\rm old}
 =K_{\rm old}(\mathsf C+\mathsf D)+K_{\rm old}\mathsf C\mathsf D:
\]
the first-order part $\mathsf C+\mathsf D$ is skew, and the symmetric part
$\mathsf C\mathsf D$ is the kernel of the two disjoint cross
pairs $(x_1,y_0)$ and $(x_0,y_1)$, spanning only two of the four
pairs of $\{x_0,x_1\}\times\{y_0,y_1\}$ (a nonstar two-edge
kernel), of normalized norm $O(1/p)$; the other orientation, $(1+\mathsf C)(1-\mathsf D)$, would make the
first-order change symmetric; this is not the mirror pairing of
\cref{lem:mirror}, which is the $(1+\mathsf C)(1+\mathsf D)$
orientation above. Within nonnegative
combinations of the accepting sign patterns a zero symmetric leading error
forces this pairing, with skew leading norm $\sqrt2\,r$. The fold record of the
construction allows signed weights: it is the signed mixture
\[
 \bigotimes_{p\in\Blk}P_{++,p}+\frac\tlaw2\bigotimes_{p\in\Blk}P_{-+,p}
 -\frac\tlaw2\bigotimes_{p\in\Blk}P_{--,p},
\]
taken after the products over the primes of the block, of three nonnegative
accepting laws, each repaired to the actual half marginal of the reference
record by an accepting Markov coupling
(\cref{def:record:mixture,lem:signed-components}). Its symmetric leading
change is still zero, its skew leading change is a three-arm star of norm
about $\sqrt{1+\tlaw^2+(1-\tlaw)^2}\,r=1.2732\,r$ at $\tlaw=127/500$, and its
variation $1+\tlaw$ does not grow with the number of primes. In the complete
sign model a zero symmetric leading error forces a one-parameter family with
this skew norm, least at $\tlaw=1/2$, whose laws have variation at least
$1+\tlaw$ for $\tlaw\ge0$; the three-component law attains this variation, and
the complete closing chooses $\tlaw=127/500$ (\cref{rem:obs:sign-frontier}).
The comparison theorem allows such a law,
because the cancellation of the kernel is applied to each nonnegative product
component and the variation multiplies only the constant of a vanishing
remainder (\cref{def:cp4}). Mirror pairing is the case $\tlaw=0$. The even
products of skew changes, which survive in $\cosh A$, are not positive
semidefinite either; positive envelopes in the algebra generated by the features
depending on a single label (the singleton features), averaged
over the reflections of the slot array in the first chamber
coordinate, price them well below the coefficient
$B$ that norm envelopes would give
(\cref{thm:envelopes,lem:oct24:mask-schur,lem:reflection-averaging}).

\paragraph{Including one fold diagonal in the dominating law.}
The bracket carries one copy of the fold diagonal,
$\Nbr=\lambda+F_{\rm fold}+H$ (\cref{lem:record:gram}), where
$F_{\rm fold}$ is the sum of the two one-layer fold copies, $H$ is
the analogous diagonal copy of the two-layer matchings, $E_{\rm add}$
below is a further pair of one-layer matching copies, $G$ is the
polarization correction (two merge records), and
$\Nbrz=\mu+E_{\rm add}+G+H$ is the auxiliary dominator of
\cref{def:record:frame}. Then the part $\Nbrz-\mu$ of the bracket
beyond the reference is invariant
under the chamber group of the first coordinate except for the polarization
correction $G$, a pair of
merge records, and the envelopes use the invariant part as the
positive background measure; its assigned price falls short of
$U_0$ by only about $2.4\cdot10^{-5}$
(\cref{lem:oct24:invariant-background}).
The price is a doubled protection of the fold records in the completion,
$v_{\rm D}\ge2w_{\rm D}k_{\rm D}$ (\cref{lem:front12:reference}(iv)), which
raises the prices of the first stage. It is affordable because the
first-stage error is a series in the two first-stage record
differences (the tags of (CP2)(b)), each tag being the difference
whose largest cut norm is attained at its own cut, whose sharp
bounds multiply there:
$R_{\rm err}<2.0445\cdot10^{-4}$ and $L_0<4.7634\cdot10^{-4}$, against about
$3.3\cdot10^{-3}$ and $7.6\cdot10^{-3}$ with one sharp factor per term
(\cref{cor:two-cut-series}).

\paragraph{The two residue classes modulo $4$.}
At $p\equiv1\pmod4$ the strict-square relation is symmetric. Labels drawn
conditionally independently from a random square clique of $\Z/p^{2\mcl}\Z$
(any two of its elements differ by an even power of $p$ times a unit square)
give a reference that is reflection positive at every cut, since a chamber cut
form is the square of a conditional mean, and that accepts every query unless
its two labels coincide (\cref{lem:frontier:iidlocal}). This event has
probability $2^{-\mcl}$ and is deleted, at a total price below
$8\cdot10^{-69}$ (\cref{lem:frontier:iidenergy}); so these primes cost
nothing at first order, and they leave the small modulus. At $p\equiv3\pmod4$
the relation is the Paley tournament, so a law that draws the two ends of a
query exchangeably cannot accept it: these primes need the oriented graph laws
and the bracket above, and the completion uses that the two-label middle law is
uniform on distinct pairs, which holds at $p\equiv3\pmod4$ and fails at $p=13$
(\cref{rem:middle-class-three}). So the laws that the $216$-slot design used for
both classes became a square-clique prefix below $2^{512}$ for class one and
the middle and upper laws for class three
(\cref{sec:frontier:iidprefix,subsec:obs:class-one}). Class three also gives
exact identities used throughout: the density ratios of the rows, the Paley
identities $\mathsf C^*=-\mathsf C$, $\mathsf C\mathsf C^*=pI-\Jall$ of the
channels of the fold records (\cref{lem:square-shift}(c)), and the exact
contraction of the index (\cref{subsec:walk:analytic}).

\subsection{The crude bound exceeds $1$: five sharper inequalities}
\label{subsec:walk:accounting}

\paragraph{Problem.}
Bounding every simultaneous change by its norm, without sharper
family bounds, without retaining the common mean, and with
$T=S+A$, gives only the crude sufficient condition
\[
 L_0+B\bigl(e^S\cosh A-1\bigr)+Be^{S+A}\Bigl(\frac{R_I}2\Bigr)^6+W<1,
\]
which is a different (stronger) requirement than $\Theta<1$: it has
no factor $63/64$, no savings $D$, and no mean term. At the baseline
values its left-hand side is about $1.603>1$
(\cref{rem:record:plain}), so this crude route cannot prove the
theorem. Five separate sharper inequalities below replace it by
$\Theta<1$.

\paragraph{Construction: five sharper inequalities.}
Each credit below bounds a separately named term or coefficient
in the closing quantities, and the five terms are distinct, so
the five savings may be added without subtracting any quantity
twice.
\begin{enumerate}[label=(\alph*)]
\item \emph{Variance credits} (\cref{sec:variance}). A repaired record has the
 actual half marginals of the record it replaces, so their difference
 annihilates the constants, and its form against the bracket is a variance:
 the energy of the diagonal copy at $\iota$ of its actual half
 marginal minus that of the product of its actual
 half marginals. Lower bounds for the latter, with positive backgrounds of the
 outside bracket, for one to four blocks and for the late pieces jointly,
 are subtracted from the reserved cost; finally $D_{\rm var}>0.0099585$.
\item \emph{Positive joint envelopes} (\cref{thm:envelopes,cor:envelope-costs}).
 The even skew products are priced below their trivial coefficient $B$: by
 positive matrices built from the single-label and star feature
 functions and averaged over the reflections of the slot array.
 The part invariant under all reflections is priced separately at
 a smaller price, since an error depending on one changed
 coordinate can first be averaged over its orbit under the
 reflections; leading and remainder prices are kept apart. With (a)
 they form the discount $D>0.46685$, and
 $\mathcal E=B(e^S\cosh A-1)-D<0.76897$, against about $1.2358$ before the
 discount, of which about $1.2016$ is the even skew part.
\item \emph{Bounding old and new records jointly} (the direct Gram
 matrix, \cref{subsec:record:direct_gram}). The matching
 envelopes of the fold records are built from a joint $2\times2$ Gram
 estimate of the old and the new record at $\iota$, instead of their norm
 envelopes. For the signed law its ideal coefficient is about $2.3842$
 ($\sqrt6$ for pure pairing); each diagonal side coefficient of the joint
 Gram bound is the larger of the signed coefficient and a
 certified auxiliary coefficient bound for that side, and only its excess over the fold diagonal
 $v_{\rm D}$ that the bracket already carries is paid
 (\cref{prop:signed-sides}). The matching exponent is $T<0.60760$ instead of
 $S+A\approx0.984$, and $J<0.25177$ instead of about $0.3668$. At the
 October~18 design this step alone took the closing from $1.1219$ to
 $1.0122$.
\item \emph{The common mean} (\cref{lem:record:common-mean}(a)). The component
 of empty support of the six signed indices is kept exactly, so the reserved
 cost is counted $63$ times out of the $64$ copies of the constant term: the
 coefficient of $L_0+\mathcal E$ is $63/64$. It saves about $0.0120$.
\item \emph{The mean credit} (\cref{prop:record:mean-credit}). Positive means of
 records of the reference, and of the chosen side diagonals that
 the (CP6) matching dominators already include, at degrees up to
 four and on the late pieces jointly, give
 $m>1.58994$ in the term $(1-m)/64$, which saves about $0.0092$.
\end{enumerate}
The result is
\[
 \Theta=\tfrac{63}{64}\bigl(L_0+\mathcal E\bigr)+J+W+\tfrac{1-m}{64}
 \approx0.757418+0.251761+0.000038-0.009218,
\]
certified as $\Theta<\PushTheta$ (\cref{sec:closing}). Without (d) and (e)
it would be about $1.0212$: in the baseline every listed credit is needed, and
the margin $1-\Theta$ is about $2.7\cdot10^{-11}$. Every credit is an exact
relation minus upper bounds for the same quantities (\cref{def:cp8}); an upper
bound is never used as a lower bound, and a certified joint bound is never
replaced by a sum of looser ones (\cref{subsec:obs:pitfalls}).

\subsection{Six parameters of the analytic stage and the estimates they control}\label{subsec:walk:analytic}

Six further parameters, from the analytic stage rather than the
local laws, were also tightened; each removes a loss that would
otherwise cost the final exponent.

\paragraph{The density exponent $\theta=1$.}
The induction bounds $Y$ by $C\Lsc^{-a}\alpha^\theta$, and the exponent of the
theorem is $a/(72-\theta)$. The analytic remainders of the cancellation and of
the large products were first bounded absolutely, by $C\Lsc^{-\sigma}$; against
$Y>C_0\Lsc^{-a}\alpha^\theta$ this needs $\alpha$ not too small, which forced
$\theta<1$ and a sparse side condition $as/(s-\theta)<\sigma$,
with $s<1$ the exponent of a bound for sparse sets. In the
six-query design of exponent $7.057\cdot10^{-6}$, with
$\theta=29/50$, this condition was within one per cent of
binding. Moment bounds that are linear in the density make every remainder
$O(\Lsc^{-\sigma}\alpha)$ (\cref{lem:record:density-can,lem:record:density-lp}),
so the induction keeps $\theta=1$ under the single condition
$a<\sigma$ (\cref{thm:record:lineardensity}); the baseline
$a\approx8.884\cdot10^{-4}$ is below
$\sigma=1/825\approx1.212\cdot10^{-3}$, and all analytic side
conditions hold on the whole
range $0<a\le9/10000$, independently of the thresholds $Q_3$ and $\Pthree$
(\cref{cert:record:analytic-stage}).

\paragraph{The Fourier cutoff $N^{1/20}$.}
The error exponent is limited by the cutoff $N^\beta$ of the lift. When the
cutoff was raised, the comparison was pointwise in the selected labels and
saved $\gamma\beta/2-\tau-72\delta$, and the large-product estimate saved
$\tau/4-\ell_0$, with $\gamma=255/1024$ the mixing exponent at the tail primes
and $\delta,\ell_0$ the moment losses (\cref{def:record:analytic-stage}); so
$\sigma<\min(\gamma\beta/2-\tau,\tau/4)\le\gamma\beta/10$, and at the cutoff
$\beta=1/36$ of the earlier versions this bound is about $6.9\cdot10^{-4}$,
below the final $a$. The stage inequalities that the certificate checks keep
this form ((S1) and (S4) of \cref{cert:record:analytic-stage}). A single
averaged comparison needs only $\gamma\beta/2-72\delta>\sigma$;
summing the raw errors over the small selections in the improved
induction uses the additional $\tau$ allowance in (S1), as proved in
\cref{prop:density:production}. The
actual-square kernel for Fourier cutoffs up to $N^{1/20}$ and prime-power
moduli (\cref{app:kernel20}) allows $\sigma=1/825$; the prime powers are needed
by the square-clique prefix.

\paragraph{The least square step.}
Fixing the labels at a set of atoms fixes a residue class modulo an integer
$D$, which is split into progressions whose step must be a square divisible by
$D$. The step is the least such square, $\sq(D)$ \eqref{eq:square-step}, where
the earlier versions used the square of the index part of $D$. The
single-level law therefore improved the step for its always-revealed modulus
$24$ from $576$ to $144$. The selected finite-depth law instead averages
$\sq(\prod_p p^{J_p})$ over its conductor depths, with the same least-square
principle at every level (\cref{prop:imp:padic-price}).

\paragraph{A fixed number of intervals.}
Earlier versions took $\Kint=2^{64}$ intervals, making the wrap negligible.
Every proof holds for arbitrary fixed $\Kint\ge2$
(\cref{rem:fixed-parameters}). The matching term grows like $\Kint^{a/6}$,
whereas the wrap decreases with $\Kint$. The selected fixed value
$\Kint=\PushK$ gives $W<3.8345\cdot10^{-5}$. Feasibility at this integer is
all the proof requires; the convexity calculation in
\cref{rem:index-convexity} guides its choice.

\paragraph{Square shifts at every conductor depth.}
At the first odd-prime level the exact square-shift contraction is
$\sqrt{p+1}/(p-1)$, optimal among doubly stochastic matrices on the
strict-square transitions (\cref{lem:square-shift}(b)). This single-level
optimality does not rule out a better distribution across conductor depths.
The selected law conditions a uniform root to have valuation less than $m$
in each prime-power coordinate, and then squares it. Its depth envelopes and finite moment are proved in
\cref{lem:imp:padic-envelope,prop:imp:padic-price}; the certified value is
$R_I<1.247139907407$ at the baseline exponent. The selected final closing
uses $R_I<1.252937308029$ (\cref{sec:four-routes-closing}). The old prime-moment
certificate remains an auxiliary majorant for this evaluation.

\paragraph{The digit pairs of the prefix.}
The square clique uses $\mcl=8192$ digit pairs. Its exception is priced below
$2^{513-\mcl+1024\mcl a}$ (\cref{lem:frontier:iidenergy}); with the earlier
$\mcl=2048$ this bound would exceed one as soon as $a>1535/2^{21}$, that is
$\eps>1.0309\cdot10^{-5}$, below the final exponent. With $8192$ pairs the
bound stays below the allowance $10^{-12}$ on the whole analytic range
$a\le9/10000$ (\cref{rem:record:prefix-walls}).

\subsection{From the local laws to the final closing}\label{subsec:walk:glance}

The examples above explain the baseline prices. The same local
interfaces support the main theorem, but its final assembly must
recompute their assignments together. \Cref{tab:walk:glance} gives
the reading path; the baseline numerical table is
\cref{tab:closing-values}, and the selected numerical ledger is in
\cref{sec:four-routes-closing}.

\begin{table}[htbp]
\centering\small
\caption{The construction's interfaces and their final assembly.}
\label{tab:walk:glance}
\begin{tabular}{@{}>{\raggedright\arraybackslash}p{0.43\textwidth}%
>{\raggedright\arraybackslash}p{0.52\textwidth}@{}}
\toprule
Output & Proof and use\\
\midrule
$\alpha^{72}\le Y$ & Reflection positivity, \cref{sec:positivity}\\
First-stage price $L_0$, dominator $B$ & Middle reference and frame,
\cref{sec:middle-reference,sec:positive-roots,sec:middle-frame}\\
Comparison prices $S,A$ & Middle and upper accepting laws,
\cref{sec:middle-records,sec:upper-comparison}\\
Discount $D$ & Fixed-marginal variance and positive envelopes,
\cref{sec:variance,sec:envelopes}\\
Matching price $T$ and joint mean & Chosen sides and disjoint mean
allocations, \cref{sec:matching-mean,subsec:joint-mean}\\
Conductor moment $R_I$ & Finite index law and summable tail,
\cref{sec:tail-index}\\
$Y\le C\Lsc^{-a}\alpha^\vartheta$ & Prefix moments and density induction,
\cref{sec:density:moments-prefix,sec:density-upgrade}\\
$\Theta<1$ & Normalized, compatible prices and the joint scalar,
\cref{sec:four-routes-closing}\\
\bottomrule
\end{tabular}
\end{table}

Each arrow from a local improvement to this closing requires its
actual marginals, completion weights, restoration bounds and
normalizing factors. A smaller isolated operator constant does not
by itself lower $\Theta$. The final assembly keeps those dependencies
explicit and evaluates every affected term at the selected parameters.
